% !TEX root = ../main.tex

\begin{abstract}[zh,title={摘\quad 要}]
  具身智能的发展正推动移动机器人从结构化实验环境走向真实场景。机器人需在地图先验不足、环境外观变化明显且人工不便持续干预的条件下，依据视觉目标自主到达指定位置。与轮式机器人相比，四足机器人具有更强的地形适应性和空间通过性，但步态运动也会引入相机抖动、速度响应非线性等系统问题。目标掩码扩散导航模型（Navigation with Goal Masked Diffusion, NoMaD）将动作扩散策略引入视觉导航，可生成多条局部航点候选轨迹；但现有验证多集中在轮式平台或离线数据集，面向四足机器人真实部署的研究仍不充分。

  针对上述问题，本文以 NoMaD 为高层视觉导航核心，以云深处 Lite3 四足机器人为目标平台，围绕``基于动作扩散策略的四足机器人视觉目标导航''展开研究，构建了从统一推理模块构建与优化、仿真闭环验证到 Jetson Orin 真机部署的完整链路。

  首先，本文以 NoMaD 为基础梳理其视觉编码器、距离预测器、目标掩码机制和扩散策略头，建立统一离线评估流程。围绕边缘设备实时推理需求，进一步构建统一视觉目标导航推理模块，涵盖去噪扩散隐式模型（Denoising Diffusion Implicit Model, DDIM）采样加速、无分类器引导（Classifier-Free Guidance, CFG）条件增强、测试时扩展（Test-Time Scaling, TTS）候选筛选和视觉编码器替换。实验结果表明，DDIM-2 能够显著降低采样开销，CFG 更适合作为场景相关可调参数，TTS-8 能够以较低代价改善候选选择，视觉编码器替换为后续真实相机适配提供了可评测框架。

  其次，本文构建由统一推理模块、有限状态机、导航主机、中层控制映射和平台接口组成的分层闭环系统，将高层航点转换为 Lite3 可执行速度指令，并屏蔽仿真与真机后端差异。基于该架构，本文在 MuJoCo 中搭建 Lite3 视觉导航闭环环境，完成目标导航、探索和多目标任务测试，验证了推理模块候选配置能够进入系统级闭环。

  此外，本文形成了面向 Jetson Orin 与 Lite3 的真机部署方案，将系统组织为实时推理链路和分层控制链路，并设计 Orin 独立调试、Lite3 联合低速验证和真实拓扑地图导航实验的分阶段联调流程。系统引入多级安全机制，使生成式视觉导航策略能够在真实平台上以可监护、可恢复的方式运行。

  实验结果表明，DDIM-2 相对基线去噪扩散概率模型（Denoising Diffusion Probabilistic Model, DDPM）实现约 4.84 倍推理加速；MuJoCo 闭环实验中，系统在多类场景下均能完成目标到达任务；Orin 端独立调试显示高层推理模块可稳定运行在约 7.0\,Hz。真机多场景目标导航实验中，基线 DDPM、最快方案（DDIM-2、CFG=0、TTS-8）、增强方案（DDIM-2、CFG=2、TTS-8）和对比方案（DDIM-3、CFG=0、无 TTS）均达到 100\% 成功率，闭环频率分别为 2.89\,Hz、6.80\,Hz、3.43\,Hz 和 5.90\,Hz。探索实验验证了无目标模式的真机可用性，室外零样本实验则揭示了强光照、复杂纹理和开阔空间对视觉编码器泛化能力的影响。综上，本文为四足机器人在弱先验环境中的视觉目标导航提供了可复现的系统框架和实验依据。
\end{abstract}

\begin{abstract}[en,title=ABSTRACT]
  The development of embodied intelligence is driving mobile robots from structured laboratory environments toward real-world scenarios. Robots need to reach visually specified goals under weak map priors, changing appearances, and limited continuous human intervention. Compared with wheeled robots, quadruped robots offer stronger terrain adaptability and mobility, while legged locomotion also introduces system-level issues such as camera disturbance and nonlinear velocity response. Navigation with Goal Masked Diffusion (NoMaD) introduces action diffusion policies into visual navigation and can generate multiple local navigation-point candidates; however, existing validations are still mainly conducted on wheeled platforms or offline datasets, and real deployment on quadruped robots remains insufficiently studied.

  This thesis takes NoMaD as the high-level navigation core and the DeepRobotics Lite3 quadruped robot as the target platform, studying visual goal navigation for quadruped robots based on action diffusion policy. A complete pipeline is built from unified inference-module construction and optimization to MuJoCo simulation and Jetson Orin real-robot deployment.

  First, the thesis analyzes NoMaD's visual encoder, distance predictor, goal mask mechanism, and diffusion policy head, establishing a unified offline evaluation workflow. To meet edge-device real-time inference requirements, a unified visual goal navigation inference module is built, covering Denoising Diffusion Implicit Model (DDIM) sampling acceleration, Classifier-Free Guidance (CFG), Test-Time Scaling (TTS) candidate selection, and visual encoder replacement. Results show that DDIM-2 reduces sampling cost, CFG is better treated as a scene-dependent tuning parameter, TTS-8 improves candidate selection with low overhead, and encoder replacement provides an evaluable framework for real-camera adaptation.

  Second, a hierarchical closed-loop system is constructed, consisting of a unified inference module, a finite-state machine, a navigation host, a mid-level control mapping module, and a platform interface. The system converts high-level navigation points into executable Lite3 velocity commands and hides backend differences between simulation and real robots. Based on this architecture, a Lite3 visual navigation environment is built in MuJoCo, where goal navigation, exploration, and multi-goal tasks are verified.

  Third, a Jetson Orin--Lite3 deployment scheme is developed. The real-robot system is organized into a real-time inference chain and a hierarchical control chain, with staged commissioning from Orin standalone debugging to low-speed Lite3 verification and real topological-map navigation. Multi-level safety mechanisms are introduced to support monitorable and recoverable real-robot execution.

  Experiments show that DDIM-2 achieves approximately 4.84$\times$ inference speedup over the Denoising Diffusion Probabilistic Model (DDPM) baseline. MuJoCo tests confirm that the system can reach targets across multiple scenario categories, and Orin standalone tests show stable high-level inference at about 7.0\,Hz. In real multi-scene navigation, the baseline DDPM setting, the fastest setting (DDIM-2, CFG=0, TTS-8), the enhanced setting (DDIM-2, CFG=2, TTS-8), and the comparison setting (DDIM-3, CFG=0, without TTS) all achieve 100\% success rates, with closed-loop frequencies of 2.89\,Hz, 6.80\,Hz, 3.43\,Hz, and 5.90\,Hz. Exploration experiments verify the usability of the goal-free mode, and outdoor zero-shot tests reveal the influence of strong illumination, complex textures, and open spaces on visual-encoder generalization. Overall, this thesis provides a reproducible system framework and experimental basis for quadruped visual goal navigation in weak-prior environments.
\end{abstract}
